%Document Class
\documentclass[a4paper,12pt,reqno]{amsart}


  \headheight0.6in
  \headsep22pt
  \textheight23cm
  \topmargin-1.7cm
  \oddsidemargin 0.5cm
  \evensidemargin0.5cm
  \textwidth15.3cm

%Packages 
\usepackage{amscd,amssymb,amsmath,amsfonts,amsthm, mathrsfs,xspace}
\usepackage{graphicx}
\usepackage{verbatim,enumerate,paralist}
\usepackage{textcomp}
%\usepackage[authoryear]{natbib}
\usepackage[pagebackref,colorlinks]{hyperref}
\hypersetup{%
  urlcolor=blue,linkcolor=blue,citecolor=blue
}
\usepackage{pdfpages}
%\usepackage{pdfsync}%  remove this when finalised

\usepackage[T1]{fontenc}
\usepackage[all,2cell,poly,knot,tips]{xy}
\UseAllTwocells
\CompileMatrices

\def\blue{\color{blue}}
\def\red{\color{red}}
\def\black{\color{black}}

%\usepackage{xypdf} % may not be necessary
\setlength{\marginparwidth}{2cm}
\newcommand{\missing}[1]{\marginpar{\color{red}missing:\\#1}}

%Theoretic Environment Setup
\newcounter{Definitioncount}
\setcounter{Definitioncount}{0}

\newtheorem*{Theorem}{Theorem}% no counter
\newtheorem{theorem}{Theorem}[section] % with counter
\newtheorem{lemma}[theorem]{Lemma}
\newtheorem{proposition}[theorem]{Proposition}
% \newtheorem*{Proposition}{Proposition}
\newtheorem{corollary}[theorem]{Corollary}
% \newtheorem*{Corollary}{Corollary}% no counter

\newtheorem*{Conjecture}{Conjecture}% no counter
\newtheorem*{Lemma}{Lemma}% no counter

\theoremstyle{definition}
\newtheorem*{Definition}{Definition}% no counter
\newtheorem*{Observation}{Observation}% no counter
\newtheorem*{Observations}{Observations}% no counter
\newtheorem*{Remark}{Remark}% no counter
\newtheorem{remark}[theorem]{Remark}
\newtheorem*{Example}{Example}% no counter
\newtheorem{example}[theorem]{Example}

\renewcommand{\theexample}{\arabic{example}}

\newtheoremstyle{fact}{\bigskipamount}{\medskipamount}{\upshape}{}{\itshape}{. }{ }{Fact}
\theoremstyle{fact}
\newtheorem*{Fact}{Fact}% no counter

\newtheoremstyle{genquest}{\bigskipamount}{\medskipamount}{\upshape}{}{\itshape}{. }{ }{General Question}
\theoremstyle{genquest}
\newtheorem*{GenQuest}{General Question}% no counter

% use \newtheoremstyle to underline the header
\newtheoremstyle{step}{2\bigskipamount}{\medskipamount}{\upshape}{}{\itshape}{. }{ }{\underline{Step~\thestep}}
\theoremstyle{step}
\newtheorem{step}{Step}[subsection]
\renewcommand{\thestep}{\arabic{step}}

\newcommand{\lra}{\longrightarrow}
\newcommand{\lla}{\longleftarrow}
\newcommand{\Lra}{\Longrightarrow}
\newcommand{\Lla}{\Longleftarrow}
\newcommand{\ldual}[1]{\mathord{{\let\nolimits\relax\sideset{^\wedge}{}{#1}}}}
\newcommand{\laction}[2]{\mathord{{\let\nolimits\relax\sideset{^{#1}}{}{#2}}}}
\newcommand{\conj}[2]{\mathord{{\let\nolimits\relax\sideset{^{#1}}{}{#2}}}}
\newcommand{\xylongto}[2][2pt]{{\UseTips{}\xymatrix @C+#1{ \ar[r]^{#2}&}}}
\newcommand{\ola}{\overleftarrow}
\newcommand{\ora}{\overrightarrow}

\DeclareMathOperator{\im}{im}

%% Script Shortcuts 
\def\CA{{\mathscr A}}
\def\CB{{\mathscr B}}
\def\CC{{\mathscr C}}
\def\CD{{\mathscr D}}
\def\CE{{\mathscr E}}
\def\CF{{\mathscr F}}
\def\CG{{\mathscr G}}
\def\CH{{\mathscr H}}
\def\CI{{\mathscr I}}
\def\CJ{{\mathscr J}}
\def\CK{{\mathscr K}}
\def\CL{{\mathscr L}}
\def\CM{{\mathscr M}}
\def\CN{{\mathscr N}}
\def\CO{{\mathscr O}}
\def\CP{{\mathscr P}}
\def\CQ{{\mathscr Q}}
\def\CR{{\mathscr R}}
\def\CS{{\mathscr S}}
\def\CT{{\mathscr T}}
\def\CU{{\mathscr U}}
\def\CV{{\mathscr V}}
\def\CW{{\mathscr W}}
\def\CX{{\mathscr X}}
\def\CY{{\mathscr Y}}
\def\CZ{{\mathscr Z}}
\def\ordm{{\mathbf{m}}}
\def\ordn{{\mathbf{n}}}
\def\ordk{{\mathbf{k}}}

\DeclareMathOperator{\Lan}{Lan}
%\DeclareMathOperator{\ev}{ev}
\DeclareMathOperator{\Skew}{Skew}
\DeclareMathOperator{\Mnd}{Mnd}
\DeclareMathOperator{\LB}{LBrMon}
\DeclareMathOperator{\Span}{Span}

\newcommand{\K}{\ensuremath{\mathscr K}\xspace}
\newcommand{\M}{\ensuremath{\mathscr M}\xspace}

\renewcommand{\L}{\ensuremath{\LB(\M)}\xspace}

%\newcommand{\Span}{\textnormal{Span}\xspace}

\newcommand{\ox}{\otimes}
\newcommand{\x}{\times}

\newcommand{\op}{\ensuremath{^{\textnormal{op}}}}


\def\theequation{{\thesection.\arabic{equation}}}


\begin{document}

\title{Orientals reinstated}

\author{Stephen Lack}
\address{Department of Mathematics, Macquarie University, NSW 2109, Australia}
\email{steve.lack@mq.edu.au}

\author{Ross Street}
\address{Department of Mathematics, Macquarie University, NSW 2109, Australia}
\email{ross.street@mq.edu.au}

\thanks{Both authors gratefully acknowledge the support of the Australian Research Council Discovery Grant DP1094883; Lack acknowledges with equal gratitude the support of an Australian Research Council Future Fellowship}

\date{\today}
\maketitle

\section{Introduction}

Monoidal categories \cite{CWM} are important in the large as bases for enriched category theory \cite{KellyBook} 
and other purposes \cite{dlVP}.  
They are also important in the small as mathematical structures themselves: in their abstract form they are pseudomonoids
(or monoidales) in monoidal bicategories \cite{{DaySt1997}, {DMcCS}}. 
It is this small viewpoint which drove the generalization of monoidal category to skew-monoidal category \cite{Szl2012}.
In particular, the relevant mathematical structure there is bialgebroid.
We \cite{{SkMon}, {SkEH}} have developed that viewpoint somewhat in the setting of quantum categories.
However, we also have found \cite{scc} that enriched category theory develops naturally with a skew-monoidal base.

At the outset \cite{{Rice}, {Kelly1964}}, partly as a test for correctness of axioms, lay the question of coherence in monoidal categories. 
That is, which diagrams involving the structural isomorphisms (constraints) commute? 
A deeper understanding \cite{aac} of coherence questions for categorical structures was that they involved knowledge 
of the free structures. 
In some interesting cases (called ``clubbable''), it was seen that the free structures on any category could essentially be 
constructed from the free structure on a single generating object (that is, on the one-object discrete category $\mathbf{1}$).
Skew-monoidal categories fall under this classification.

The notion of free needs to be made clear. One meaning is that we consider the underlying 2-functor from the 
2-category of structures, with morphisms preserving structure up to coherent isomorphism, to the 2-category of categories;
the left biadjoint to this gives the free structure on each category. These free structures are only unique up to equivalence.
In the case of monoidal categories, we can take the free structure on $\mathbf{1}$ simply to be the discrete category $\mathbb{N}$
with the natural numbers as objects and addition as the tensor product. One might call this notion of free {\em bicategorical}.   

Another meaning is that we consider the underlying 2-functor from the 2-category of structures, with morphisms {\em strictly}
preserving the structure, to the 2-category of categories; the left adjoint to this gives the free structure on each category. 
These free structures are unique up to isomorphism. This is freeness in the {\em strict} sense. A discussion of these matters can be found in \cite{btc}.

The present paper describes a model for the free skew-monoidal category $\mathrm{Fsk}$ on $\mathbf{1}$ in the strict sense.
Our construction involves the orientals of \cite{aos} and the Tamari lattices of \cite{HT}.

\section{Axioms, non-axioms and an important example}

A {\em (left) skew-monoidal category} is a category $\CC$ equipped with an object $I$ (called the {\em unit} or {\em skew unit}),
a functor $\otimes : \CC \times \CC \lra \CC$ (called {\em tensor product}), and
natural families of {\em lax constraints} having the directions 
\begin{equation}
\alpha_{XYZ} : (X\otimes Y)\otimes Z \lra X\otimes (Y\otimes Z)
\end{equation}
\begin{equation}
\lambda_X : I\otimes X \lra X
\end{equation}
\begin{equation}
\rho_X :  X \lra X\otimes I
\end{equation}
subject to five conditions: 

\begin{equation}\label{pentagon}
\xymatrix{
& (W\otimes X)\otimes (Y\otimes Z) \ar[rd]^-{{\alpha_{W,X,Y\otimes Z}}}  & \\
((W\otimes X)\otimes Y)\otimes Z \ar[ru]^-{{\alpha_{W\otimes X,Y,Z}}} \ar[d]_-{{\alpha_{W,X,Y}} \otimes 1_Z} & & W\otimes (X\otimes (Y\otimes Z))  \\
(W\otimes (X\otimes Y))\otimes Z \ar[rr]_-{{\alpha_{W,X\otimes Y,Z}}} & & W\otimes ((X\otimes Y)\otimes Z) \ar[u]_-{1_W\otimes {\alpha_{X, Y,Z}}} }
\end{equation}

\begin{equation}\label{leftunit}
\xymatrix{
(I\otimes X)\otimes Y \ar[rd]_{\lambda_{X} \otimes 1_Y}\ar[rr]^{\alpha_{I,X,Y}}   && I\otimes (X\otimes Y) \ar[ld]^{\lambda_{X\otimes Y}} \\
& X\otimes Y  &
}
\end{equation}

\begin{equation}\label{midunit}
\xymatrix{
(X\otimes I)\otimes Y \ar[rr]^-{\alpha_{X,I,Y}} && X\otimes (I\otimes Y) \ar[d]^-{1_X\otimes \lambda_Y} \\
X\otimes Y \ar[u]^-{\rho_X \otimes 1_Y} \ar[rr]_-{1_{X\otimes Y}} && X\otimes Y}
\end{equation}

\begin{equation}\label{rightunit}
\xymatrix{
(X\otimes Y)\otimes I \ar[rr]^{\alpha_{I,X,Y}}   && X\otimes (Y\otimes I)  \\
& X\otimes Y \ar[lu]^{\rho_{X \otimes Y}}  \ar[ru]_{1_X\otimes \rho_Y}&
}
\end{equation}

\begin{equation}\label{unitunit}
\xymatrix{
I \ar[rd]_{\rho_I}\ar[rr]^{1_I}   && I  \\
& I\otimes I \ar[ru]_{\lambda_I} & .
}
\end{equation}

Notice that we obtain idempotents
$$\epsilon^{\ell}_{X,Y} : (X\otimes I) \otimes Y \stackrel{\alpha_{X,I,Y}} \lra X\otimes (I \otimes Y)
\stackrel{1_X \otimes \lambda_X} \lra X\otimes Y
\stackrel{\rho_X \otimes 1_Y} \lra (X\otimes I)\otimes Y \ ,$$
$$\epsilon^{r}_{X,Y} : X\otimes (I \otimes Y) \stackrel{1_X \otimes \lambda_X} \lra X\otimes Y
\stackrel{\rho_X \otimes 1_Y} \lra (X\otimes I)\otimes Y \stackrel{\alpha_{X,I,Y}} \lra X\otimes (I \otimes Y)$$
and
$$\epsilon_0 : I\otimes I \stackrel{\lambda_I} \lra I \stackrel{\rho_I}\lra I\otimes I $$
which are not identities in general. Moreover, 
$$\alpha_{X,I,Y} : ((X\otimes I) \otimes Y, \epsilon^{\ell}_{X,Y}) \lra (X\otimes (I \otimes Y), \epsilon^{r}_{X,Y})$$
is a morphism of idempotents. 

We shall give an example of a skew-monoidal category that plays an important role in this paper.
Along with it, we fix some notation.
Recall the algebraists' simplicial category $\mathbf{\Delta}$.
The objects are the finite ordinals $$\mathbf{n} = \{0,1,\dots ,n-1\} \ .$$
The morphisms are order-preserving functions $\xi : \mathbf{m} \lra \mathbf{n}$. 
For each such function, we put
$$\xi_{\ell} = \{i\in \mathbf{m-1} : \xi (i) = \xi (i+1)\}$$
and
$$\xi_{r} = \{j \in \mathbf{n} : j \notin \mathrm{im} \xi \} \ .$$
Notice that $\xi_{\ell} =\varnothing$ means $\xi$ is injective and $\xi_{r} =\varnothing$ means $\xi$ is surjective. 
Then there is a unique factorization $\xi = \partial \circ \sigma$ where 
$\sigma_{\ell} = \xi_{\ell}$, $\sigma_{r} = \varnothing$, $\partial_{\ell} = \varnothing$, and $\partial_{r} = \xi_r$.  

It is well known \cite{osed} that $\mathbf{\Delta}$ is the free strict monoidal category containing a monoid.
The unit object is $\mathbf{0}$ and tensor product is ordinal sum $\mathbf{m}\otimes \mathbf{n} = \mathbf{m+n}$, $\xi\otimes \zeta = \xi \mathbf{+} \zeta$.
The generic monoid is $\mathbf{1}$ with the only multiplication and unit possible.

Let $\mathbf{\Delta}_{\bot}$ denote the subcategory of $\mathbf{\Delta}$ consisting of non-empty finite ordinals $\mathbf{n}$
and first-element-preserving functions; that is, the $\xi$ with $0\notin \xi_r$.
The tensor product on 
$\mathbf{\Delta}$ restricts to $\mathbf{\Delta}_{\bot}$ 
but we replace the unit by the terminal object $\mathbf{1}$.
This gives a skew-monoidal category with identity associativity constraint, where 
$$\lambda_{\mathbf{n}} : \mathbf{1+n} \lra \mathbf{n}$$ is the surjection with 
$(\lambda_{\mathbf{n}})_{\ell} = \{0\}$, and 
$$\rho_{\mathbf{n}} : \mathbf{n} \lra \mathbf{n+1}$$ is the injection with
$(\rho_{\mathbf{n}})_{r} = \{n\}$.

When we write $\mathbf{\Delta}_{\bot}$ in future, we mean it equipped with the above skew-monoidal structure. 

Each morphism $\xi : \mathbf{m}\lra \mathbf{n}$ in $\mathbf{\Delta}_{\bot}$
has a right adjoint $\xi^* : \mathbf{n}\lra \mathbf{m}$ as order-preserving functions. 
The formula is
\begin{equation}\label{radj}
\xi^*(j) = \mathrm{max}\{i : \xi(i) \le j \} \ .
\end{equation}

Let $\mathbf{\Delta}_{\top}$ denote the subcategory of $\mathbf{\Delta}$ consisting of non-empty finite ordinals $\mathbf{n}$
and last-element-preserving functions.
This becomes a right skew-monoidal category in the obvious way; in fact, in a dual way as reinforced by the next result.

Each morphism $\xi : \mathbf{m}\lra \mathbf{n}$ in $\mathbf{\Delta}_{\top}$
has a left adjoint $\xi^! : \mathbf{n}\lra \mathbf{m}$ as order-preserving functions. 
The formula is
\begin{equation}\label{ladj}
\xi^!(j) = \mathrm{min}\{i : j \le \xi(i) \} \ .
\end{equation}

\begin{proposition}\label{bottop} 
Formula \eqref{radj} defines an isomorphism of categories
$$(\Delta_{\bot})^{\mathrm{op}} \cong \Delta_{\top} \ .$$ 
If $\zeta : \mathbf{m}\lra \mathbf{n}$ in $\mathbf{\Delta}_{\top}$ is
the image of $\xi : \mathbf{n}\lra \mathbf{m}$ in $\mathbf{\Delta}_{\bot}$ under the isomorphism then 
$$ \zeta_r = \xi_{\ell}  \ \text{  and  } \  \zeta_{\ell} = \{i-1 : i\in \xi_r \}  \ .$$ 
\end{proposition}
\begin{proof}
A functor between finite ordinals preserves limits if and only if it preserves last element.
A functor between finite ordinals preserves colimits if and only if it preserves first element.
So the isomorphism follows from the adjoint functor theorem and the fact that the right adjoint of a functor is the left Kan extension of the identity along the functor. 
The second sentence follows using composition of adjunctions $\tau \dashv \sigma \dashv \partial$
where $\partial$ and $\tau$ are injective, $\sigma$ is surjective, $\partial_r = \sigma_{\ell} = \{i\}$ and $\tau_r = \{i+1\}$
for some natural number $i$.      
\end{proof}

\section{Lax monoidal categories}

A {\em lax-monoidal category} is a category $\CC$ equipped with an object $I$ (called the {\em unit} or {\em lax unit}),
a functor $\otimes : \CC \times \CC \lra \CC$ (called {\em tensor product}), and
natural families of {\em lax constraints} having the directions 
\begin{equation}
\alpha_{XYZ} : (X\otimes Y)\otimes Z \lra X\otimes (Y\otimes Z)
\end{equation}
\begin{equation}
\lambda_X : I\otimes X \lra X
\end{equation}
\begin{equation}
\rho_X :  X\otimes I \lra X
\end{equation}
subject to five conditions: 

\begin{equation}\label{pentagon}
\xymatrix{
& (W\otimes X)\otimes (Y\otimes Z) \ar[rd]^-{{\alpha_{W,X,Y\otimes Z}}}  & \\
((W\otimes X)\otimes Y)\otimes Z \ar[ru]^-{{\alpha_{W\otimes X,Y,Z}}} \ar[d]_-{{\alpha_{W,X,Y}} \otimes 1_Z} & & W\otimes (X\otimes (Y\otimes Z))  \\
(W\otimes (X\otimes Y))\otimes Z \ar[rr]_-{{\alpha_{W,X\otimes Y,Z}}} & & W\otimes ((X\otimes Y)\otimes Z) \ar[u]_-{1_W\otimes {\alpha_{X, Y,Z}}} }
\end{equation}

\begin{equation}\label{leftunit}
\xymatrix{
(I\otimes X)\otimes Y \ar[rd]_{\lambda_{X} \otimes 1_Y}\ar[rr]^{\alpha_{I,X,Y}}   && I\otimes (X\otimes Y) \ar[ld]^{\lambda_{X\otimes Y}} \\
& X\otimes Y  &
}
\end{equation}

\begin{equation}\label{midunit}
\xymatrix{
(X\otimes I)\otimes Y \ar[d]_-{\rho_X \otimes 1_Y} \ar[rr]^-{\alpha_{X,I,Y}} && X\otimes (I\otimes Y) \ar[d]^-{1_X\otimes \lambda_Y} \\
X\otimes Y  \ar[rr]_-{1_{X\otimes Y}} && X\otimes Y}
\end{equation}

\begin{equation}\label{rightunit}
\xymatrix{
(X\otimes Y) \ar[rd]_{\rho_{X \otimes Y}}\otimes I \ar[rr]^{\alpha_{I,X,Y}}   && X\otimes (Y\otimes I)  \ar[ld]^{1_X\otimes \rho_Y}  \\
& X\otimes Y &
}
\end{equation}

\begin{equation}\label{unitunit}
\xymatrix{
I\otimes I \ar[rr]^{\lambda_I = \rho_I}   & & I \ .
}
\end{equation}

\begin{example}\label{Dbt}
Let $\Delta_{\bot \top}$ denote the subcategory of the algebraic simplicial category 
$\Delta$ with those objects 
$\ordm$ with $1\le m$ and morphisms those which preserve top and bottom element.  
Actually, $\Delta_{\bot \top}\cong \Delta^{\mathrm{op}}$. 
The tensor product of $\Delta$, which is ordinal sum, restricts to $\Delta_{\bot \top}$.
As lax unit object for $\Delta_{\bot \top}$ we take $\mathbf{1}$.
We have morphisms 
$\lambda_{\mathbf{m}} = \sigma_0 : \mathbf{1+m} \lra \mathbf{m}$ and $\rho_{\mathbf{m}} = \sigma_{\mathbf{m-1}} : \mathbf{m +1} \lra \mathbf{m}$
which are the order-preserving surjections equating $0$ and $1$, and equating 
$m-1$ and $m$, respectively.
The axioms for a lax-monoidal category are easily checked.   
\end{example}

\section{Bracketing functions}\label{Bf}

A {\em left bracketing function} (lbf) is a function $\ell:\ordm \lra \ordm$ satisfying
\begin{itemize}
\item[(i)] $\ell(j) \le j$ for all $j\in \ordm$;
\item[(ii)] $\ell(j)\le i \le j<m$ implies $\ell(j)\le \ell(i)$;
\item[(iii)] $\ell(m-1)=m-1$.
\end{itemize}
Some consequences are:
\begin{itemize}
\item[(iv)] $\ell(0) =0$;
\item[(v)] $\ell$ is idempotent.
\end{itemize}
As with any functions into a linearly ordered set, the lbf are ordered using value-wise order in $\ordm$.
This is the Tamari lattice $\mathrm{Tam}_m$; see \cite{{Tamari51},{Tamari62}, {HT}} .

A {\em right bracketing function} (rbf) is a function $r:\ordm \lra \ordm$ satisfying
\begin{itemize}
\item[(i)] $i \le r(i)$ for all $i\in \ordm$;
\item[(ii)] $0\le i\le j \le r(i)$ implies $r(j)\le r(i)$;
\item[(iii)] $r(0)=0$.
\end{itemize}
Some consequences are:
\begin{itemize}
\item[(iv)] $r(m-1) =m-1$;
\item[(v)] $r$ is idempotent.
\end{itemize}

\begin{proposition}\label{bijbf}
The following equations determine a bijection between lbfs $\ell:\ordm \lra \ordm$ 
and rbfs $r:\ordm \lra \ordm$: for $0<i,j<m-1$,
$$r(i) = \mathrm{min} \{j : \ell(j) < i \le j \}$$
$$\ell(j) = \mathrm{max} \{i : i \le j < r(i) \} .$$
\end{proposition}
\begin{proof} 
Given an lbf $\ell$, let $r$ be defined as in the proposition 
together with the requirement that it should preserve $0$ and $m-1$. 
We need to see first that $r$ is an $rbf$.
As every $j$ with $\ell(j) < i \le j$ has $i \le j$, we obtain $i \le r(i)$.  
To prove property (ii) for $r$, suppose $i\le j\le r(i)$.
Then $\ell(k)<i \le k$ implies $j\le k$. If $r(i) < r(j)$ then there exists
$k$ with $\ell(k)< i\le k$ but not $\ell(k)< j\le k$. Since we know $j\le k$,
we must have $j \le \ell(k)$. Then we have the contradiction
$i\le j\le \ell(k) < i$. So $r(j) \le r(i)$ as required.

Now we need to show that we can recover $\ell$ from the second formula in the
proposition. We need to see that $\ell(j)$ is the biggest $i$ such that $i\le j$ and, if
$\ell(k) < i \le k$, then $j<k$. To see that $i = \ell(j)$ does have the property notice that
$\ell(j) \le j$ by (i) for $\ell$, and, by (ii) for $\ell$, $\ell(k) < \ell(j) \le k$ and 
$k\le j$ would imply $\ell(j)\le \ell(k)$, a contradiction. 
To see that $\ell(j)$ is the biggest such $i$, suppose we had
$\ell(j) < i$ with $i\le j$ and $\ell(k) < i \le k$ implying $j<k$. Then $j$ is such a $k$.
So $j<j$, a contradiction.

This proves half of the proposition. The other half is by left-right symmetry.  
\end{proof}

When $\ell$ and $r$ correspond under the bijection of Proposition~\ref{bijbf}
we put $\ora{\ell} = r$ and $\ell=\ola{r}$.

We emphasise that $\ell(0)=0=r(0)$, $\ell(m-1)=m-1=r(m-1)$ and, for $0<i,j<m-1$,
we have $\ell(j)< m-1$ and $r(i)>0$. 
Yet we can have $\ell(j)=0$ 
(when $\{i : i \le j < r(i) \}=\varnothing$) or $r(i)=m-1$ 
(when $\{j : \ell(j) < i \le j \}=\varnothing$).

\begin{proposition}\label{ordbf}
The bijection of Proposition~\ref{bijbf} is order preserving.
\end{proposition}
\begin{proof}
Assume $\ell \le \ell_1$. Then $\{ j : \ell_1 (j)< i\le j \} \subseteq \{j: \ell (j)< i \le j\} $.
So $\ora{\ell} \le \ora{\ell_1}$.
\end{proof}

Consequently, we may consider the elements of 
$\mathrm{Tam}_m$ to be rbfs rather than lbfs.

\begin{proposition}\label{fixbf}
Suppose $\ell$ and $r$ correspond under the bijection of Proposition~\ref{bijbf}.
\begin{itemize}
\item[(i)] For $0<i<m-1$, $\ell(i)=i$ if and only if $r(i) \neq i$.
\item[(ii)] For $0<i\le j <m-1$, if $\ell (j) = i$ then $r(i) \neq j$.
\item[(iii)] If $0<i$ and $r(i)<m-1$ then $\ell r(i) \le \ell(i-1)$.
\item[(iv)] If $j<m-1$ and $0<\ell(j)$ then $r(j+1) \le r\ell(j)$. 
\item[(v)] If $0<i$, $0<\ell(i-1)$ and $r(i)<m-1$ then $\ell r(i) \neq \ell(i-1)$
implies $r(i) = r\ell(i-1)$.  
\end{itemize}
\end{proposition}
\begin{proof}
For (i) we have $\ell(i)\neq i$ iff $\ell(i)<i$ iff $\ell(i)<i\le i$ iff $\ora{\ell}(i)\le i$
iff $\ora{\ell}(i) = i$.   

For (ii), if $i=\ell(j)$ then $j$ cannot be a $k$ with $\ell(k)<i \le k$ (let alone the first such).
So $\ora{\ell}(i) \neq j$. 

For (iii), since $r(i)<m-1$, $r(i)$ is a $k$ satisfying $\ell(k) < i\le k$ so $\ell r(i)<i$;
so $\ell r(i)\le i-1$. But, by property (ii) for a lbf, 
$\ell r(i)\le i-1< i \le r(i)$ implies $\ell r(i)\le \ell (i-1)$. 

Then (iv) follows from (iii) by left-right symmetry. 

Now to (vi). Assume $0<i$, $0<\ell(i-1)$, $r(i)<m-1$ and $\ell r(i) \neq \ell(i-1)$.
By (ii), $\ell r(i) < \ell(i-1)$. Now $r\ell(i-1)$ is the minimum $j$ with 
$\ell(j)<\ell(i-1)\le j$. Yet $\ell r(i) < \ell(i-1) \le r(i)$. So $r\ell(i-1)\le r(i)$.
But, by (iv), $r(i)\le r\ell(i-1)$. So $r(i)= r\ell(i-1)$. 
\end{proof}

\begin{proposition}\label{lbf}
Suppose $\xi : \ordm \lra \ordn$ preserves order and top and bottom element
so that there are adjoints $\xi^! \dashv \xi \dashv \xi^*$. 
If $\ell:\ordm \lra \ordm$ is an lbf then so is $\xi \ell \xi^*:\ordn \lra \ordn$. 
If $r:\ordm \lra \ordm$ is an rbf then so is $\xi r \xi^!:\ordn \lra \ordn$.
\end{proposition}
\begin{proof} 
Clearly $\ell_1=\xi \ell \xi^*$ preserves top element since all three factors do. Also $\ell_1(j) = \xi \ell \xi^*(j) \le \xi \xi^* (j) \le j$ using that (i) holds for $\ell$ and $\xi$ preserves order. 
It remains to prove (ii) for $\ell_1$.
Now $\ell_1(j) \le i \le j$ implies $\ell\xi^*(j) \le \xi^*(i) \le \xi^*(j)$ which implies, using (ii) for $\ell$, $\ell\xi^*(j) \le \ell\xi^*(i)$. 
However $\xi$ is order preserving, so $\ell_1(j)\le \ell_1(i)$.
The result for $r$ is obtained left-right symmetrically.
\end{proof}

\section{Motivation, orientals and the Tamari lattice}

Our goal is to discover a nice model for the free skew-monoidal category $\mathrm{Fsk}$ generated by a single object $X$. 
The objects are clearly words in the letters $I$ and $X$, bracketed meaningfully pairwise. 
For example, there is the object 
\begin{equation}
(((X(IX))(XX))X) \ ;
\end{equation}
but we can write this as a triangulated polygon
\begin{equation}\label{heptagon1}
\xymatrix{
& 0 \ar[rr]{} \ar[rrrd]{} \ar[rddd]{}  \ar[ld]_{X}  & & 6  &  \\
1 \ar[rrdd]{} \ar[d]_{I} & & & & 5 \ar[lu]_{X} \\
2  \ar[rrd]_{X} & & & & 4  \ar[u]_{X} \\
& &  3 \ar[rruu]{}  \ar[rru]_{X} &
}
\end{equation}
where the edges $i\lra i+1$ are labelled by the symbols $I$ or $X$.
We can omit the labels $I$ since anything not labelled by $X$ can be understood to be labelled by $I$.

We shall use the convention that each subset $x=\{x_0,x_1, \dots ,x_r\}$ of $\mathbf{m}\in \mathbf{\Delta}$ organizes itself as a list 
$x=(x_0x_1\dots x_r)$ with $x_0<x_1< \dots <x_r$. 
We may omit the curly brackets around some sets of such subsets and may also omit the commas between elements. 
For $0\le p \le r$, the {\em $p$-th face} $x\partial_p$ of 
$x$ is obtained by deleting $x_p$ from $x$.
The face is called {\em odd} or {\em even} according as $p$ is odd or even.
We write $x^-$ for the set of odd faces of $x$ and we write $x^+$ for the set of even faces. 
Because we start with an even $0$, there is one more even than odd face.
If $H$ is a set of subsets of $\ordm$ of cardinality $r+1$, we define
\begin{eqnarray}
H^- = \{x^- : x\in H\} \ \text{  and  } \ H^+ = \{x^+ : x\in H\} \ . 
\end{eqnarray}
Also, we define
\begin{eqnarray}
H^{\mp} = H^- \smallsetminus H^+ \ \text{  and  } \ H^{\pm} = H^+ \smallsetminus H^- \ . 
\end{eqnarray}

For each $\xi : \mathbf{m} \lra \mathbf{n}$ in $\mathbf{\Delta}$, we write
$\xi x$ for $\{\xi(x_0)\xi(x_1) \dots \xi(x_r)\} \subseteq \mathbf{n}$. 
So for each $\xi : \mathbf{m} \lra \mathbf{n}$ in $\mathbf{\Delta}_{\bot}$
and $y=(y_0y_1\dots y_s) \subseteq \mathbf{n}$, we have
$$\xi^* y = \{\xi^*(y_0)\xi^*(y_1) \dots \xi^*(y_s)\} \subseteq \mathbf{m} \ .$$ 

Now we recall the {\em orientals} $\CO_m$ \cite{aos}. This is the free $m$-category on the $m$-simplex.
However, we need no more than the underlying 4-category structure. 
The objects (= 0-cells) of $\CO_m$ are the natural numbers $p$ with $0 \le p \le m$. 
A morphism (= 1-cell) $a : p \lra q$ implies $p\le q$ and $a=(a_0 a_1\dots a_r)$ is a subset of $\mathbf{q}\smallsetminus \mathbf{p}$
with $a_0 =p$; we think of $a$ as the path 
$$p\stackrel{(pa_1)}\lra a_1\stackrel{(a_1a_2)}\lra a_2\lra \dots \stackrel{(a_rq)}\lra q \ .$$
According to \cite{aos} we think of the morphism $a : p \lra q$ as a set of doublets 
$$a=(pa_1)(a_1a_2)\dots (a_rq) \ .$$  
For our purposes, we are interested in the morphism $b_m : 0 \lra m$ which is $0\stackrel{(0m)}\lra m$ and
the morphism $e_m : 0 \lra m$ which is the path $0\stackrel{(01)}\lra 1 \stackrel{(12)}\lra 2 \dots  \stackrel{(\overline{m-1}m)}\lra m$. 

Let $\CT_m = \CO_m (0,m)(b_m,e_m)$ as a 2-category. We shall examine this more explicitly. 

The objects of $\CT_m$ are 2-cells $S:b_m \Lra e_m$ in $\CO_m$.
A description of these, adapted from \cite{aos}, is as follows: 
$S$ is a set of three-element subsets $x=(x_0x_1x_2)$ of $\mathbf{m+1}$ 
satisfying the conditions:
\begin{itemize}
\item[(a)] for $x,y\in S$, if two of the three equations $x_0=y_0$, $x_1=y_1$, $x_2=y_2$ hold then $x=y$;
\item[(b)] for $x,y\in S$, if $x_1=y_0$ then $x_2\neq y_1$;
\item[(c)] $b_m = S^{\mp}$;
\item[(d)] $e_m = S^{\pm}$.
\end{itemize} 
Some consequences are:
\begin{itemize}
\item[(e)] the function $S\lra \{1,2,\dots ,m-1\}$, taking $x$ to its middle element $x_1$, is a bijection;
\item[(f)] for $x,y\in S$, if $x_0<y_1\le x_1$ then $x_0\le y_0$;
\item[(g)] for $x,y\in S$, $y_2\le x_1$ if and only if $x_0< y_1<x_1$.
\end{itemize} 
Injectivity of the function in (e) follows from the alternating position (AP) condition proved in \cite{aos}; surjectivity is a simple induction related to the excision of extremals
algorithm in \cite{aos}. 

These objects $S$ are in bijection with triangulations of the polygon with $m$ sides.
For example, the triangulation \eqref{heptagon1} of the heptagon has 
\begin{equation}
S=(013)(123)(035)(345)(056) \ .
\end{equation} 
It will be useful to have some notation involving the inverse $t_S$ of the bijection in (e) above:
\begin{eqnarray}\label{tlr}
t_S : \{1,2,\dots ,m-1\} \lra S \nonumber \\
t_S(i) = (\ell_S(i-1), i , r_S(i)+1) \ .
\end{eqnarray}
Notice that, in the process, we are defining functions 
\begin{eqnarray}
\ell_S : \mathbf{m-1} \lra \mathbf{m-1}
\end{eqnarray}
and 
\begin{eqnarray}\label{rightfns}
r_S: \{1,2,\dots ,m-1\} \lra \{1,2,\dots ,m-1\} \ .
\end{eqnarray}
However, by putting $\ell_S(m-1)=m-1$ and $r_S(0)=0$, we extend them to bottom-and-top-preserving functions
\begin{equation}\label{lr}
\ell_S, r_S : \mathbf{m} \lra \mathbf{m} \ .
\end{equation}

We put 
\begin{equation}
S_{\ell} = \{i : 0\le i<m-1, \ell_S(i)=i \}
\end{equation}
and
\begin{equation}
S_r = \{i :  0<i\le m-1, r_S(i)=i \} \ .
\end{equation}

\begin{proposition}\label{Sbijlbf} 
The formula
$$S=\{(\ell(i-1), i , \ora{\ell}(i)+1) : 0<i<m-1\}$$
establishes a bijection between lbf morphisms $\ell : \mathbf{m} \lra \mathbf{m}$ 
and  2-cells $S:b_m \Lra e_m$ in $\CO_m$. 
\end{proposition}
\begin{proof}
Suppose $\ell$ is a lbf. Put $r=\ora{\ell}$. 
Define $S$ as in the proposition. 
We need to prove properties (a), (b), (c), (d) for $S$.
Property (e) is clear. So only one case for property (a) remains. 
Take $x=(\ell(i-1), i , r(i)+1)$
and $y=(\ell(j-1), j , r(j)+1)$ in $S$. We need to see that 
$\ell (i-1) = \ell (j-1)$ and $r(i)=r(j)$ imply $i=j$.
Assume $i<j$ since the hypotheses are symmetric.
From $r(i)=r(j)$, we see that any $k$ with
$\ell (k) < i \le k$ will have $r (j) \le k$. However, $j-1$ is just such
a $k$ since $\ell (j-1) = \ell (i-1) \le i-1 < i \le j-1$. So $r (j) \le j-1$.
But then $j \le r(j) \le j-1$ gives a contradiction.   

For property (b), take $x$ and $y$ in $S$ as above. 
This time we need to see that $x_1=y_0$ implies $x_2\neq y_1$;
that is, $i=\ell(j-1)$ implies $r(i)+1\neq j$. We do have
$j-1<m-1$ and $0<i=\ell(j-1)$, so, using (iv) of Proposition~\ref{fixbf},
we have $r(i)=r\ell (j-1) \ge r(j) \ge j > j-1$. So $r(i)+1\neq j$.  

Now we look at property (c). Since, by (i), $\ell(i-1) < i$, there must exist an 
$i$ with $\ell (i-1) = 0$.
Choose the largest such $i$. Then $\ell(k)<i\le k$ implies $\ell(k) \le i-1 <k$,
so, by (ii) for an lbf, $\ell(k) \le \ell(i-1)=0$. So $\ell(k)=0$. 
By maximality of $i$, $k\le i-1$, a contradiction. 
So the set of $k$, for which $r(i)$ is the minimum, is empty.
So $r(i) =m-1$. This shows that $b_m\subseteq S^-$.
Yet $b_m\cap S^+=\varnothing$ since there is no $x\in S$
with $x_0<0$ or $x_2>m$. 
So $b_m\subseteq S^{\mp}$.
Next we need to show that, for any $x\in S$, if $x^-\neq (0m)$ then $x^- \in S^+$.
Assume $x=(\ell(i-1), i , r(i)+1)$ and $0< \ell(i-1)$ (since the case
$r(i)+1<m$ will follow by left-right symmetry). 
We need to show that there exists either $y$ or $z$ in $S$ with
$y\partial_0 = (l(i-1),r(i)+1)$ or $z\partial_2 = (l(i-1),r(i)+1)$.
The last two conditions force $y=(\ell(j-1),j, r(j)+1)$ to have $j=\ell(i-1)$
and $z=(\ell(k-1),k, r(k)+1)$ to have $k=r(i)+1$. So we must see that,
with these choices for $j$ and $k$, either  $r(j)=r(i)$ or $\ell(k-1)=\ell(i-1)$.
Assume $\ell(k-1)\neq \ell(i-1)$; that is, $\ell r (i) \neq \ell(i-1)$. 
Then part (v) of Proposition~\ref{fixbf} implies $r(i)=r\ell(i-1)=r(j)$.

We also need to show $e_m = S^{\pm}$. Take $(i,i+1)\in e_m$. 
By part (i) of Proposition~\ref{fixbf}, either $i = r(i)$ or $\ell(i) =i$. 
So either $(i,i+1) = (\ell(i-1), i,r(i)+1)\partial_0$
or $(i,i+1) = (\ell(i), i+1,r(i+1)+1)\partial_2$. 
However we cannot have $(i,i+1) = (\ell(j-1), j,r(j)+1)\partial_1$
since $\ell(j-1)$ and $r(j)+1)$ have $j$ strictly between them 
and so cannot be consecutive.
So $e_m\subseteq S^{\pm}$. 
Now we prove $S^{\pm}\subseteq e_m$ by showing that $x\in S$
with $(x_1,x_2)\notin e_m$ or $(x_0,x_1)\notin e_m$ implies $x^+ \subseteq S^-$.
Suppose $x=(\ell(i-1), i , r(i)+1)$ and $(i , r(i)+1)\notin e_m$.
Then $i<r(i)$, so, by part (i) of Proposition~\ref{fixbf}, we have $\ell(i)=i$. 
Let $j$ be the largest such that $i\le j\le r(i)$ and $\ell(j)=i$.
By part (v) of Proposition~\ref{fixbf}, we have $r(j) = r(i)$.
So $x\partial_0 = (i , r(i)+1) = (\ell(j-1),j,r(j)+1)\partial_1 \in S^-$. 
The case $x\partial_2\notin e_m$ is similar. 

So far, we have defined a function $\mathrm{Tam}_m \lra \CT_m$ which
is obviously injective. If we accept that the two sets have the same cardinality
(given by a Catalan number), the bijection follows. However, we will
describe the inverse function and complete the proof. 

Suppose $S:b_m \Lra e_m$ is a 2-cell. Using property (e) for $S$, 
define top-and-bottom-preserving functions 
$\ell = \ell_S$ and $r=r_S$ by \eqref{tlr} and \eqref{lr}. 
Then $\ell(i-1)< i < r(i)+1$ for $0<i<m-1$.
This gives property (i) for $\ell$ to be an lbf  (and for $r$ to be an rbf).
To prove property (ii) for $\ell$, take $\ell(j)\le i\le j<m$ and put
$x=(\ell(j),j+1,r(j+1)+1)$ and $y=(\ell(i),i+1,r(i+1)+1)$ which are in $S$.
Then $x_0<y_1\le x_1$. By property (f) for $S$, $x_0\le y_0$; so $\ell(j)\le \ell(i)$. 
It remains to show that $r = \ora{\ell}$. 
This breaks into two parts. First we must see that $r(i)<m-1$ implies 
$\ell r(i)<i\le r(i)$. For this, take $x=(\ell r(i),r(i)+1,r(r(i)+1)+1)$ and
$y=(\ell (i-1),i,r(i)+1)$ so that $x_1=y_2$. By the `only if' part of property (g) for $S$,
$x_0<y_1<y_2$, our desired result. The second part is minimality; that is, to show
$\ell(j)<i\le j$ implies $r(i)\le j$. Put $x=(\ell r(i),r(i)+1,r(r(i)+1)+1)$ and
$y=(\ell(i-1),i,r(i)+1)$ so that $x_0<y_1<x_1$. 
By the `if' part of property (g) for $S$, $y_2\le x_1$; that is, $r(i)\le j$.
\end{proof}

Now let us look at the morphisms of $\CT_m$.
As 3-cells in $\CO_m$, they are well defined in \cite{aos} as sets $\theta$ of
four-element subsets $x=(x_0x_1x_2x_3)$ of $\mathbf{m+1}$ satisfying
well formedness and movement conditions; also see \cite{{pc},{pc:c}}.
The indecomposable $\theta : S\lra T$ are precisely the singleton sets $\{ x\}$
such that
\begin{equation}\label{xmove}
T = (S\smallsetminus x^{-})\cup x^{+} \ ,
\end{equation}
where 
\begin{eqnarray*}
x^{-} = (x_0x_1x_2)(x_0x_2x_3) \ \text{ and  }  \ x^{+} = (x_0x_1x_3)(x_1x_2x_3) \ .
\end{eqnarray*}
So $x^{-}$ consists of the odd faces of $x$ and $x^{+}$ the
even faces in the sense of \cite{aos}, although the notation is from \cite{pc}. 
Let us put
\begin{equation}
xS=T \ \text{   and   } \ S=Tx
\end{equation}
when \eqref{xmove} holds.
In pictures, this means that $S$ should contain the domain of
$x=(x_0x_1x_2x_3)$ in the diagram \eqref{indec3}
as part of its triangulation and $T$ should be obtained from $S$ by 
replacing that part by the codomain.

\begin{eqnarray}\label{indec3}
\xymatrix{
x_0 \ar[r]^-{} \ar[rd]^-{} \ar[d]_-{} & x_3 \\
x_1 \ar[r]_-{}  & x_2 \ar[u]^-{} }
\qquad
\xymatrix{
\stackrel{x} \lra}
\qquad
\xymatrix{
x_0 \ar[r]^-{} \ar[d]_-{} & x_3 \\
x_1 \ar[r]_-{} \ar[ru]^-{}  & x_2 \ar[u]^-{} }
\end{eqnarray}
In terms of bracketings, it means that $T$ is obtained from $S$ by moving just 
one set of brackets to the right:
$$((W_0((W_1W_2)W_3))W_4) \lra ((W_0(W_1(W_2W_3)))W_4)$$
or
$$(W_0(((W_1W_2)W_3)W_4)) \lra (W_0((W_1(W_2W_3))W_4))$$
where the $W_i$ are meaningfully binarily bracketed words in $I$ and $X$.

Now let us look at the 2-cells of $\CT_m$.
As 4-cells in $\CO_n$, they are well defined in \cite{aos} as sets $a$ of
five-element subsets $x=(x_0x_1x_2x_3x_4)$ of $\mathbf{m+1}$ satisfying
well formedness and movement conditions.
The indecomposable $a : \theta \Lra \phi$ are precisely 
the singleton sets $\{ x\}$ such that
\begin{equation}
\phi = (\theta \smallsetminus x^{-})\cup x^{+} \ ,
\end{equation}
where 
$$x^{-} = (x_0x_1x_2x_4)(x_0x_2x_3x_4)$$ and 
$$x^{+} = (x_1x_2x_3x_4)(x_0x_1x_3x_4)(x_0x_1x_2x_3) \ .$$
In diagram~\eqref{O4}, we see the only non-identity 2-cell in $\CT_4$.

\begin{equation}\label{O4}
\xymatrix{
(012)(023)(034) \ar[d]_{(0234)}^(0.8){\phantom{AAAAAAAA}}="1" \ar[rr]^{(0123)}  && (013)(123)(034) \ar[d]^{(0134)}_(0.8){\phantom{AAAAAAAA}}="2" \ar@{=>}"1";"2"^-{(01234)}
\\
(012)(024)(234) \ar[rd]_-{(0124)} && (014)(123)(134) \ar[ld]^-{(1234)} 
\\
& (014)(124)(234) 
}
\end{equation}

The following result is immediate on combining Proposition~\ref{Sbijlbf}, \cite{HT} and \cite{aos}.

\begin{proposition}\label{TLatt} For 2-cells $S,T:b_n\Lra e_n : 0\lra n$ in the $n$-category 
$\CO_n$, the following conditions are equivalent:
\begin{itemize}
\item[(i)] $\ell_S \le \ell_T$;
\item[(ii)] $r_S \le r_T$;
\item[(iii)] there exists a 3-cell $S\lra T$ in $\CO_n$;
\item[(iv)] for all $x\in S$ and $y\in T$, if $x_1=y_1$ then $x_0\le y_0$; 
\item[(v)] for all $x\in S$ and $y\in T$, if $x_1=y_1$ then $x_2\le y_2$;
\item[(vi)] for all $x\in S$ and $y\in T$, if $x_1=y_1$ then $x_0\le y_0$ and $x_2\le y_2$.
\end{itemize}
\end{proposition}

So this means that the Tamari lattice $\mathrm{Tam}_m$ 
(see Section~\ref{Bf}) is obtained from the 2-category
$\CT_m$ by identifying all the 2-cells. That is, there is a 2-functor
$\CT_m \lra \mathrm{Tam}_m$ (where the only 2-cells in 
$\mathrm{Tam}_m$ are the identities); and every 2-functor 
$\CT_m \lra \CC$ which takes all 2-cells in $\CT_m$ to identities
factors uniquely through $\CT_m \lra \mathrm{Tam}_m$. 

If $\ell_S \le \ell_T$ in $\mathrm{Tam}_n$, we define the {\em distance} 
from $S$ to $T$ by
\begin{equation}
\mathrm{d}(S,T) = \mathrm{min} \{\# \theta : \theta : S \lra T \text{ in } \CT_n \} \ .
\end{equation}
Clearly, $\mathrm{d}(S,S) = 0$ and $\mathrm{d}(S,U) \le \mathrm{d}(S,T)+\mathrm{d}(T,U)$ for $S\le T\le U$. 
As an example obtained by looking at diagram~\eqref{O4}, we have
$$\mathrm{d}((012)(023)(034), (014)(124)(234)) = 2 \ .$$
  

For $m,n > 0$, there is an order-preserving function 
\begin{equation}
\star : \mathrm{Tam}_m \times \mathrm{Tam}_n \lra \mathrm{Tam}_{m+n}
\end{equation}
defined by 
$$S\star T = S\cup \{(0,m,m+n)\} \cup (m+T) \ .$$
Here $m+T = \{(m+i,m+j,m+k) : (i,j\k) \in T\}$.  
In other words,
\begin{equation*}
  t_{S\star T} (i) =
      \begin{cases}
         t_S(i) & \text{if } \  i < m  \\
         (0,m,m+n) & \text{if } \  i = m  \\
         (m,m,m)+t_T(i-m) & \text{if } \ i>m \ .
       \end{cases}
\end{equation*}
The existence of the 3-cell 
$$(0,m,m+n,m+n+k) : (S\star T)\star U \Lra S\star (T\star U)$$
in $\CO_{m+n+k}$ shows we have lax associativity in the form
\begin{equation}\label{Tamlaxassoc}
(S\star T)\star U \le S\star (T\star U) \ .
\end{equation} 
Note that $\mathrm{Tam}_1 = \{\varnothing\}$ is the terminal poset; 
its element is the identity 2-cell of $b_1 = e_1 : 0\lra 1$ in $\CO_1$.
This element does not act as a lax identity since we have 
$$\varnothing \star S =  \{(0,1,1+m)\} \cup (1+S)$$ 
and 
$$S \star \varnothing = S\cup \{(0,m,m+1)\}$$
in $\mathrm{Tam}_{m+1}$; so these cannot be Tamari compared with 
$S\in \mathrm{Tam}_{m}$.  

\section{Shrink morphisms}

\begin{proposition}\label{surju}
Suppose $\sigma :\ordm \lra \ordn$ is a surjective order-preserving function
with right adjoint $\sigma^*$. 
Suppose $u\subseteq \ordm$ and $v\subseteq \ordn$. 
The following conditions are equivalent:
\begin{itemize}
\item[(i)] $\sigma^*$ induces a bijection between $v$ and $u$;
\item[(ii)] $\sigma$ induces a bijection between $u$ and $v$, and 
$\sigma_{\ell}\cap u = \varnothing$.
\end{itemize}
\end{proposition}
\begin{proof}
Both conditions imply $u$ and $v$ have the same cardinality: write
$u=(u_0u_1\dots u_{p})$ and $v=(v_0v_1\dots v_{p})$.
Recalling Proposition~\ref{bottop}, we have
$\sigma_{\ell}=\sigma^*_{r}$; so
$\sigma_{\ell}\cap u = \varnothing$ is equivalent to saying 
$u$ is in the image of $\sigma^*$. 
Since $\sigma$ is surjective, we also have $\sigma \sigma^*=1_{\ordn}$.
Now (i) means $u_k=\sigma^*(v_k)$ for $0\le k\le p$; so certainly this shows
$u$ is in the image of $\sigma^*$, while applying $\sigma$ gives
$\sigma(u_k) = v_k$, showing that $\sigma$ induces a bijection 
between $u$ and $v$. 

Now assume condition (ii). Then $v_k=\sigma(u_k)$. But also $u_k$
is in the image of $\sigma^*$, so $u_k=\sigma^*(i)$ for some $i$.
Then $v_k=\sigma(u_k)=\sigma \sigma^*(i) = i$ and we have
$u_k=\sigma^*(v_k)$.  
\end{proof}

\begin{Definition}
Let $\mathrm{Fsk}_{\ell}$ denote the category whose objects are a triples 
$(m,u,S)$ where $m>0$, $u\subseteq \ordm$ and $S\in \CT_m$.
A morphism $\sigma : (m,u,S) \lra (n,v,T)$,
called a {\em shrink morphism}, is a surjective
order-preserving function $\sigma : \ordm \lra \ordn$ such that
\begin{itemize}
\item[(a)] $\sigma^*$ induces a bijection between $v$ and $u$; 
\item[(b)] $\ell_T = \sigma \ell_S \sigma^*$;
\item[(c)] $r_T = \sigma r_S \sigma^!$.
\end{itemize}   
Clearly shrink morphisms compose and we do have a category.
\end{Definition}

\begin{proposition}
Suppose $\sigma :\ordm \lra \ordn$ is a surjective order-preserving function
with right adjoint $\sigma^*$ and suppose $(m,u,S)$ is an object of $\mathrm{Fsk}_{\ell}$.
There exists a shrink morphism of the form $\sigma : (m,u,S) \lra (n,v,T)$ 
if and only if
\begin{itemize}
\item[(i)] $\sigma_{\ell}\cap u = \varnothing$, and
\item[(ii)] $\ora{\sigma \ell_S \sigma^*} = \sigma r_S \sigma^!$.
\end{itemize}
\end{proposition}
\begin{proof}
The necessity of the conditions follows from
Propositions \ref{surju},\ref{bijbf},and \ref{Sbijlbf}. 
For sufficiency, notice that (i) implies the restriction of 
$\sigma$ to $u$ is injective (since $i<j$ in $u$ with $\sigma(i)=\sigma(j)$
implies $\sigma(i)=\sigma(i+1)$ contrary to $i\notin \sigma_{\ell}$).  
Thus $\sigma$ induces a bijection between $u$ and $v=\sigma u$. 
By Proposition~\ref{surju}, we have (a) of the Definition.
By (b) and (c) of the Definition, we are forced to put
$\ell_T = \sigma \ell_S \sigma^*$ and $r_T = \sigma r_S \sigma^!$.
Using (ii), this defines $T\in \CT_n$ by Propositions \ref{lbf} and \ref{Sbijlbf}.  
We have our shrink morphism as required.
\end{proof}

\begin{proposition}
If $\sigma : (m,u,S) \lra (n,v,T)$ is a shrink morphism and $T \le T^{\prime}$
(meaning $\ell_T\le \ell_{T^{\prime}}$ in $\mathrm{Tam}_n$) then there
exists a shrink morphism of the form 
$\sigma : (m,u,S^{\prime}) \lra (n,v,T^{\prime})$ 
with $S\le S^{\prime}$. 
\end{proposition}
\begin{proof}
Using Proposition~\ref{TLatt}, we only need to look at the case where
there is an indecomposible 2-cell $y=(y_0y_1y_2y_3) : T\Lra T^{\prime}$. 
Then $T^{\prime}$ is obtained from $T$ by replacing the elements
$(y_0y_1y_2)$ and $(y_0y_2y_3)$ by $(y_0y_1y_3)$ and $(y_1y_2y_3)$
with the same middle terms. So
\begin{equation*}
  \ell_{T^{\prime}} (k) =
      \begin{cases}
          \ell_T (k) & \text{if } \  k \neq y_2-1  \\
           y_1 & \text{if } \ k = y_2-1 \ .
       \end{cases}
\end{equation*}
Put $x_1=\sigma^*(y_1-1)+1$ and $x_2=\sigma^*(y_2-1)+1$ noting
that $\sigma^*$ is strictly order preserving; so $x_1<x_2$. Define
a function $\ell_{S^{\prime}}:\ordm\lra\ordm$ by
\begin{equation*}
  \ell_{S^{\prime}} (h) =
      \begin{cases}
          \ell_S (h) & \text{if } \  h \neq x_2-1  \\
           x_1 & \text{if } \ h = x_2-1 \ .
       \end{cases}
\end{equation*}
Applying adjointness to the inequality
$$\sigma \ell_S \sigma^*(y_2-1)=\ell_T(y_2-1)=y_0\le y_1-1$$
yields the inequality
$$\ell_S(x_2-1)=\ell_S\sigma^*(y_2-1)\le \sigma^*(y_1-1)=x_1-1<x_1=\ell_{S^{\prime}}(x_2-1) \ .$$
So $\ell_S\le \ell_{S^{\prime}}$ value-wise.
We claim that $\ell_{S^{\prime}}$ is an lbf. 
For property (i), we have $\ell_{S^{\prime}}(h)=\ell_S(h)\le h$ for $h\neq x_2-1$
while $\ell_{S^{\prime}}(x_2-1)=x_1\le x_2-1$. For property (ii), assume
$\ell_{S^{\prime}}(h)\le i < h <m$. If neither $i$ nor $h$ is $x_2-1$ then
$\ell_{S^{\prime}}(h)=\ell_S(h) \le \ell_S(i)=\ell_{S^{\prime}}(i)$.
If $i=x_2-1$ then $\ell_{S^{\prime}}(h) = \ell_S(h) \le x_2-1<h$; so 
$\ell_{S^{\prime}}(h) = \ell_S(h) \le \ell_S(x_2-1) < x_1 = \ell_{S^{\prime}}(i)$.  
In the case $h=x_2-1$ we need to show that $x_1\le i < x_2-1$ implies 
$x_1\le \ell_S(i)$. 
\red{I am stuck at this point and maybe this approach is wrong. However, 
I hope Steve has a proof of this proposition.}
\end{proof}

\begin{proposition}
If $\sigma : (m,u,S) \lra (n,v,T)$ and $\sigma : (m,u^{\prime},S^{\prime}) \lra (n,v,T)$ 
are a shrink morphisms then $u=u^{\prime}$ and either $S \le S^{\prime}$
or $S^{\prime} \le S$. 
\end{proposition}
\begin{proof}
\red{Help Steve!!}
\end{proof}

\section{Swell morphisms and general morphisms}

\begin{Definition}
Let $\mathrm{Fsk}_r$ denote the category whose objects are a triples 
$(m,u,S)$ where $m>0$, $u\subseteq \ordm$ and $S\in \CT_m$.
A morphism $\partial : (m,u,S) \lra (n,v,T)$,
called a {\em swell morphism}, is an injective
order-and-bottom-preserving function 
$\partial : \ordm \lra \ordn$ such that
\begin{itemize}
\item[(a)] $\partial$ induces a bijection between $u$ and $v$; 
\item[(b)] $\ell_S = \partial^* \ell_T \partial^{**}$;
\item[(c)] $r_S = \partial^* r_T \partial$.
\end{itemize}   
Clearly swell morphisms compose and we do have a category.
\end{Definition}

\begin{proposition}
There is an isomorphism of categories
$$\mathrm{Fsk}_r^{\mathrm{op}}\cong \mathrm{Fsk}_{\ell}$$
which is the identity on objects and takes $\partial$ to $\partial^*$.
\end{proposition}
\red{I am confused here. The duality seems to require it to be $\partial^{**}$
which induced a bijection between $u$ and $v$. But in the examples it is
as in (a) above.} 

\section{A bijection}

It is well known that rbfs are in bijection with elements of the Catalan monoid.
We give an account of this. In order to construct an inverse for a certain function 
$\Theta$ from the Tamari lattice to the Catalan monoid, 
we require a construction on certain partial maps.

We will define a new partial map $h^{\prime} : \mathbf{n+1} \lra \mathbf{n}$ from a given one $h :\mathbf{n+1} \lra \mathbf{n}$.
We will let $x_h \subseteq \mathbf{n+1}$ denote the domain of $h$ and use the same symbol $h$ for the function $h: x_h \lra \ordn$. 
Let $a_h$ be the first $i \in x_h$ with $h(i) > 0$, if it exists.

Obtain $x_{h^{\prime}} \subseteq \mathbf{n+1}$ from $x_h$ 
by deleting precisely $h(i)$ 
of the elements directly preceding $i$ in $x_h$. 
Then define $h^{\prime} : x_{h^{\prime}} \lra \mathbf{n}$ to agree with $h$ on all elements of $x_{h^{\prime}}$
except $a_h$ where $h^{\prime}(a_h) = 0$.

Notice that $a_h < a_{h^{\prime}}$ if the latter exists.   

The so-called {\em Catalan monoid} is defined by 
$$\mathrm{Cln}_n = \{ \text{order-preserving functions} \ g : \mathbf{n-1} \rightarrow \mathbf{n-1} : g(i) \le i \ \text{for all} \ i\in \mathbf{n-1} \} $$
with the multiplication given by composition of functions.

There is a function $$\Theta : \mathrm{Tam}_n \lra \mathrm{Cln}_n $$
defined by $\Theta (r) (j) = \#r^{-1} \{1, 2, \dots , j\}$ for $0 \le j \le n-2$
for any rbf $r:\ordn \lra \ordn$.

\begin{proposition} $\Theta$ is a bijection.
\end{proposition}
\begin{proof}
Without going into the full proof, we define the inverse of $\Theta$.
Let $g\in \mathrm{Cln}_n$.
We recursively define a sequence of functions $h^m : x^m \lra \ordn$, $1 \le m \le q$, where $q$ is the number of non-empty fibres of $g$. 
Put $x^1 = \ordn$ and $h^1 = \# g^{-1}(i)$. 
Then $h^{m+1} = (h^m)^{\prime}$ with $x^{m} = x_{h^m}$.
Notice that $$\bigcup_{1\le m \le q} x_m \smallsetminus x_{m+1} = \{1, \dots , n-1\} \ . $$
For $i \in x_m \smallsetminus x_{m+1}$, define $$r(i) = a_{h^m}-1 \ .$$ 
Then $\Theta(r) = g$.
\end{proof}

The pointwise order on $\mathrm{Tam}_n$ obviously transports across the bijection $\Theta : \mathrm{Tam}_n \lra \mathrm{Cln}_n$ to an order on 
$\mathrm{Cln}_n$.
However, this description takes some work to disentangle. 
On the other hand, it is easy to describe the indecomposable
inequalities $g_1 < g_2$ in $\mathrm{Cln}_n$.  

\begin{proposition} The Tamari order on $\mathrm{Tam}_n$ transports across $\Theta$ to the partial order on 
$\mathrm{Cln}_n$ generated by the relation
$g_1 < g_2$ if there exists $0<i<j \le n-1$ such that 
\begin{equation*}
  g_1 (k) =
      \begin{cases}
         g_2 (k) + 1 & \text{if } \  i \le k \le j \ , \\
         g_2 (k) & \text{otherwise }.
       \end{cases}
\end{equation*}
\end{proposition}

Clearly $\mathrm{Cln}_n$ is a monoid with zero under composition of endofunctions of $\mathbf{n-1}$.
However, it is not an ordered monoid. 
To see this, take $g_0$, $g_1$, $g_2$, $g_3$ to be the functions
whose values are $012$, $011$, $002$, $001$, respectively. 
While $g_0 < g_2$, when we multiply
on the left by $g_1$ we see that $g_1 g_0 = g_1$ and $g_1 g_2 = g_3$; yet $g_1 \nleq g_3$. 


\begin{thebibliography}{00}

\bibitem{Cohn1965} P.M. Cohn, \textit{Universal Algebra} (Harper \& Row, 1965) pp. 269 and 289, Ex. 9.\label{Cohn1965}

\bibitem{DaySt1997} Brian J. Day and Ross Street, \textit{Monoidal bicategories and Hopf algebroids}, Advances in Math. \textbf{129} (1997) 99--157. \label{DaySt1997}

\bibitem{DMcCS} Brian J. Day, Paddy McCrudden and Ross Street, \textit{Dualizations and antipodes}, Applied Categorical Structures \textbf{11} (2003) 229--260. \label{DMcCS]}

\bibitem{HT} Dani\`ele Huguet and Dov Tamari, \textit{Problems of associativity: a simple proof for the lattice property of
systems ordered by a semi-associative law}, J. Comb. Theory (A) \textbf{13} (1972) 7--13.\label{HT}  

\bibitem{btc}  Andr\'e Joyal and Ross Street, \textit{Braided tensor categories}, Advances in Mathematics \textbf{102} (1993) 20--78.\label{btc}

\bibitem{Kelly1964} G. Max Kelly, \textit{On MacLane's conditions for coherence of natural associativities, commutativities, etc.}, Journal of Algebra \textbf{1} (1964) 397--402.\label{Kelly1964}

\bibitem{aac} G. Max Kelly, \textit{An abstract approach to coherence}, Lecture Notes in Math. \textbf{281} (Springer-Verlag, 1972) 106--147.\label{aac}

\bibitem{KellyBook} G. Max Kelly, \textit{Basic concepts of enriched category theory}, London Mathematical Society Lecture Note Series \textbf{64} (Cambridge University Press, Cambridge, 1982). \label{KellyBook}

\bibitem{Steve20Dec2012} Stephen Lack and Ross Street, \textit{Triangulations, orientals, and skew monoidal categories} (Working notes, 20 December 2012).\label{Steve20Dec2012} 

\bibitem{SkMon} Stephen Lack and Ross Street, \textit{Skew monoidales, skew warpings and quantum categories}, Theory and Applications of Categories \textbf{27(3)} (2012) 27--46.
\label{SkMon}

\bibitem{SkEH} Stephen Lack and Ross Street, \textit{A skew-duoidal Eckmann-Hilton argument and quantum categories}, (submitted; see http://arxiv.org/abs/1210.8192).\label{SkEH}

\bibitem{osed} F. William Lawvere, \textit{Ordinal sums and equational doctrines}, Reprints in Theory and Applications of Categories \textbf{18} (2008) 113--121.\label{osed}

\bibitem{Rice} Saunders Mac Lane, \textit{Natural associativity and commutativity}, Rice Univ. Studies \textbf{49} (1963) 28--46.\label{Rice}

\bibitem{CWM} Saunders Mac Lane, \textit{Categories for the Working Mathematician}, Graduate Texts in Mathematics \textbf{5} (Springer-Verlag, 1971).\label{CWM} 

\bibitem{aos} Ross Street, \textit{The algebra of oriented simplexes}, J. Pure Appl. Algebra \textbf{49} (1987) 283--335.\label{aos}

\bibitem{pc} Ross Street, \textit{Parity complexes}, Cahiers de topologie et g\'eom\'etrie diff\'erentielle cat\'egoriques \textbf{32} (1991) 315--343.\label{pc}

\bibitem{pc:c} Ross Street, \textit{Parity complexes: corrigenda}, Cahiers de topologie et g\'eom\'etrie diff\'erentielle cat\'egoriques \textbf{35} (1994) 359--361.\label{pc:c}

\bibitem{dlVP} Ross Street, \textit{Monoidal categories in, and linking, geometry and algebra}, Bulletin Belgian Math. Society  (to appear; see http://arxiv.org/abs/1201.2991).\label{dlVP}

\bibitem{scc} Ross Street, \textit{Skew-closed categories}, Journal of Pure and Applied Algebra (2012), doi: 10.1016/j.jpaa.2012.09.020; also see http://arxiv.org/abs/1205.6522).\label{scc}

\bibitem{Szl2012} Kornel Szlach\'anyi, \textit{Skew-monoidal categories and bialgebroids}, Advances in Mathematics \textbf{231(3--4)} (2012) 1694--1730.
%arXiv:1201.4981v1 (2012).
\label{Szl2012}

\bibitem{Tamari51} Dov Tamari, \textit{Monoides pr\'eordonn\'es et cha\^ines de Malcev}, 
Doctorat\`es-Sciences Math\'ematiques Th\`ese de Math\'ematiques (Parisa, 1951).\label{Tamari51} 

\bibitem{Tamari62} Dov Tamari, \textit{The algebra of bracketings and their enumeration}, Nieuw Archief voor Wiskunde \textbf{10} (1962) 131--146.\label{Tamari62}




\end{thebibliography}

\end{document}